\subsection{Reached sections on the actual polarized model}
\label{subsec:ftb-reached-sections}

Retain the integral uniformization and its coefficient bundles from
\S\ref{subsec:ftb-expansion-regularity}.
The following uses its period/cubical identifications, including the
quadratic valuation \(Q\); it introduces no replacement model.
Put
\[
 b_Q(x,y)=Q(x+y)-Q(x)-Q(y).
\]
The polarized period pairing gives a symmetric bilinear \(b_Q\),
\(b_Q(x,x)=2Q(x)\), and
\[
 Q(x+km)=Q(x)+k\,b_Q(m,x)+k^2Q(m).
\]

An actual generator of the principal theta line has an integral
primitive coefficient at \(0\).
The coefficient functional equation transports it from \(0\) to all
of \(X\); more generally the same equation shows that the inequalities
defining \(\nu\ge r\) can be checked on representatives of \(X/kX\).
Indeed the coefficient line-bundle translation for
\(\chi\mapsto\chi+km\) changes the scalar valuation by
\(kQ(m)+b_Q(m,\chi)\).  The displayed quadratic identity says that
this is precisely
\((Q(\chi+km)-Q(\chi))/k\).
For \(k=1\) there is only the representative \(0\).
Thus the primitive theta generator has \(\nu\ge0\).
This is the source's
\texttt{GeneratorSectionExpansionNuNonneg}, using
\texttt{ValuationMinRepresentatives} with the representative set
\(\{0\}\).

We give the two operations on this assertion explicitly.
For a CM endomorphism \(a\) with \(a^*a=aa^*=N\), the adjoint
compatibility of the period pairing gives
\[
 b_Q(a^*\chi,a^*\chi)=b_Q(\chi,N\chi),\qquad
                         Q(a^*\chi)=NQ(\chi).
\]
The normalized pullback at line powers \(k\to kN\) has coefficients
supported on \(a^*X\); the coefficient at \(a^*\chi\) is the
pullback of the coefficient at \(\chi\) through the corresponding
endomorphism of the proper abelian quotient, with the canonical
rigidified coefficient-bundle identification.
This identification has an integral generator: it is the cubical
Rosati identity on the integral quotient, normalized at its identity.
Pullback takes an integral section to an integral section.
Consequently
\[
 V_{\ge r+Q(\chi)/k}(\Theta_\chi)
 \Longrightarrow
 V_{\ge r+Q(a^*\chi)/(kN)}
                       (([a]^*\Theta)_{a^*\chi}).
\]
The other coefficients vanish and satisfy every lower bound.
Therefore normalized CM pullback preserves the inequality \(\nu\ge r\)
in the forward direction used here.
The source calls these the
\texttt{PullGenStatement}, \texttt{QAlStatement}, and
\texttt{PullbackValuationForwardBound};
the last is obtained by exactly this coefficient argument.
Its application to sections includes the source restriction
\(k=1\) or \(k\ge4\).

For a symmetric torsion translation, let \(p\) be its uniformized
point at line power \(k\), with period label \(y\in X\) and proper
abelian coordinate \(b\) satisfying \(kb=c(y)\).
Properness extends \(b\) to the valuation ring.
For each \(\chi\), compare its chosen generic Poincare generator with
an integral generator over this extension.
The scalar \(t_\chi\) in this comparison satisfies
\begin{equation}\label{eq:ftb-poincare-ratio}
                     k\,v(t_\chi)=b_Q(y,\chi).
\end{equation}
Here is the cubical valuation calculation behind this formula.
Taking \(k\) tensor powers changes the \(b\)-argument to \(kb=c(y)\).
The cubical multiplication identification compares this tensor with
the period-pairing generator at \((y,\chi)\).
All generators extending over the valuation ring have unit ratios;
their valuations are zero.  The period pairing has valuation
\(b_Q(y,\chi)\).
The extra \(kQ(y)\) term in the theta rigidification occurs on both
sides and cancels.  Thus the remaining identity of valuations is
\(kv(t_\chi)+kQ(y)=b_Q(y,\chi)+kQ(y)\), proving
\eqref{eq:ftb-poincare-ratio}.
This is the tensor calculation printed in
\texttt{ValuationThetaRatioEqBilinear.v\_ratio\_gt};
its hypotheses include the same cubical structure and period
rigidifications of the uniformization.

If \(q_w\) is the valuation of the single remaining line generator,
and the translation lift has scalar normalization \(\zeta\), its
coefficient map shifts \(\chi\) to \(\chi+y\).
By \eqref{eq:ftb-poincare-ratio} it gives, for all \(r\),
\[
 \nu(\Theta)\ge r
 \quad\Longleftrightarrow\quad
 \nu(\zeta\Phi_p\Theta)\ge r+c,\qquad
 c=q_w+v(\zeta)-Q(y)/k .
\]
The equality follows by subtracting
\(Q(\chi+y)/k\) from the coefficient valuation:
the \(\chi\)-dependent term
\(b_Q(y,\chi)/k\) cancels exactly.
This is the content of \texttt{TWMain.tw\_of}, rather than the
unsupported assertion that an arbitrary finite-order matrix preserves
an arbitrary lattice.

\begin{srlemma}[The finite-order shift argument]\label{lem:ftb-periodic-nu}
Suppose \(A\) is a finite-order operator on these coefficient
families and, for a fixed rational \(c\),
\(\nu(s)\ge r\) if and only if \(\nu(As)\ge r+c\), for every \(s,r\).
Then \(\nu(s)\ge r\) if and only if \(\nu(As)\ge r\).
\end{srlemma}
\begin{proof}
Let \(A^M=1\), with \(M\ge1\).
Iteration gives
\[
      \nu(s)\ge r\quad\Longleftrightarrow\quad
                              \nu(s)\ge r+Mc.
\]
If \(c=0\), there is nothing to prove.
Otherwise the monotone predicate \(P(r):=[\nu(s)\ge r]\) is periodic
with a nonzero rational period.  If it holds at one rational threshold,
it holds at every rational threshold: use multiples of the absolute
period to exceed any prescribed threshold, then decrease the bound.
In particular \(P(r)\) and \(P(r-c)\) are equivalent.
Combining this with the one-step shift gives the assertion.
This proof permits the zero family; no unproved finiteness or upper
bound for \(\nu\) is needed.
\end{proof}

An actual symmetric \(k\)-torsion lift acts with finite order on
coefficient families.
Its lift through the PEL torsion identification can differ from a
chosen uniformized lift only by a root of unity.
Such a scalar has valuation zero.
Lemma~\ref{lem:ftb-periodic-nu} and the preceding exact shift formula
therefore show that symmetric torsion translation preserves
\(\nu\ge0\).
The actual finite-order and root-of-unity comparisons are those of
\texttt{SymmetricLiftActsRootOfUnityOnExpansions}, used in
\texttt{SymmetricTorsionTranslateSectionNuNonneg}.

Define a \emph{reached} section inductively as follows.
Start from the actual principal theta generator at line power \(1\).
One may apply a normalized CM pullback
\(a^*a=N\), with \(N\) odd, from line power \(k\) to \(kN\),
provided \(k=1\) or \(k\ge4\).
One may also apply a normalized symmetric torsion translation at
the current line power.
This is exactly the inductive predicate \texttt{Reach}.

\begin{srproposition}[Reached sections are regular on the same chart]
\label{prop:ftb-reached-regular}
For a good integral chart with its actual PEL uniformization,
every reached section extends over that integral model.
\end{srproposition}
\begin{proof}
Base change the chart and all its cubical data to an algebraically
closed valued extension of its fraction field.
The base coefficient calculation, the normalized-pullback argument,
and the symmetric-translation argument above prove \(\nu\ge0\)
by induction on the construction of the reached section.
The pullback and translation identities commute with this base
change: pulling the corresponding line-bundle isomorphisms and their
normalized generators gives the same identities, by uniqueness of
the rigidified normalization.
Apply Proposition~\ref{prop:ftb-expansion-regularity} to descend the
nonnegative coefficient family and extend the section over the
original chart.
This is \texttt{ReachedSectionRegularOnIntegralModel.regular\_of\_reach}.
The proposition is relative to the actual integral uniformization
and its compatible expansion correspondence, as stated above.
\end{proof}

\subsection{The exact mixed-translation pullback}
\label{subsec:ftb-mixed-pullback}

\begin{srlemma}[The reached section behind a mixed translate]
\label{lem:ftb-mixed-pullback}
Let \(N\) be odd, and let
\[
 N=\beta\alpha n',\qquad
 \alpha\bar\alpha=A,\quad\beta\bar\beta=B,\qquad
                                N^2=(n')^2AB .
\]
Here the first equality is in \(\mathcal O_E\), and \(A,B,n'\) are
nonnegative integers.
Let \(s\) be the actual line-power-one theta section.
Suppose \(x_1,x_2\) satisfy
\[
 x_1^{N^2}=x_2^{N^2}=1,\quad
 ([\alpha\beta]x_1)^{(n')^2}=1,\quad
 ([\beta]x_2)^{(n')^2A}=1 .
\]
If \(F\) is the prescribed normalized mixed translate of
\([N]^*s\) at these two points, then
\[
                         F=[\beta]^*h
\]
for a reached section \(h\) at line power \((n')^2A\).
All pullbacks and translates have their fixed symmetric normalization.
\end{srlemma}
\begin{proof}
Construct, in order,
\[
 G_0=[n']^*s,\quad
 G_1=t_{[\alpha\beta]x_1}^*G_0,\quad
 K_1=[\alpha]^*G_1,\quad
 h=t_{[\beta]x_2}^*K_1 .
\]
The displayed torsion conditions make both translations defined.
Every intermediate line power is odd; \((n')^2\) is \(1\) or at
least \(4\), so these are permitted \texttt{Reach} steps.
Composition of normalized pullbacks gives
\([\alpha\beta]^*[n']^*s=[N]^*s\).
Normalized pullback commutes with symmetric translation by replacing
the translating point by its image.
Apply these two identities first to the \(x_1\)-translate and then
to the \(x_2\)-translate.
Uniqueness of normalized odd-level symmetric translates identifies
\([\beta]^*h\) with \(F\) exactly, including its scalar phase.
This proves the assertion of
\texttt{MixedTranslatePullbackEqReachedSection.uu\_eq\_pull\_reach}.
\end{proof}

For a harmonic polynomial \(P\) with row degrees \(a_{\tau,i}\),
let
\[
                 \operatorname{Nm}_a(\gamma)
                  =\prod_{\tau,i}\tau(\gamma)^{a_{\tau,i}} .
\]
Every nonzero monomial of \(P\) has total degree \(d_a\):
sum its row degrees over all rows.
Applying the differential chain rule in each eigendirection gives
\[
 \mathcal D_P([\gamma]^*h)
          =\operatorname{Nm}_a(\gamma)\,\mathcal D_P(h),\qquad
 \operatorname{Nm}_a(m)=m^{d_a}.
\]
These identities are valid for arbitrary coefficient-field
coefficients of \(P\), by linearity; the source first clears their
common denominator when applying its integral-polynomial version.

\begin{srproposition}[The mixed-translate valuation bound]
\label{prop:ftb-mixed-value-bound}
In an actual good integral chart suppose the harmonic connection
evaluation of an integral germ has valuation at least \(-\rho\) for
all integral harmonic polynomials of these row degrees and all line
powers.
Retain the hypotheses of Lemma~\ref{lem:ftb-mixed-pullback}, and suppose
\(n'\) is a unit in the chart DVR.
Then
\[
 v\left(
  \frac{\operatorname{Nm}_a(\alpha)}{N^{d_a}}
                         \mathcal D_P(F)\right)\ge-\rho .
\]
\end{srproposition}
\begin{proof}
The reached section \(h\) is regular by
Proposition~\ref{prop:ftb-reached-regular}, so its germ is an integral
power series in this chart.
The pullback chain rule and \(N=\beta\alpha n'\) give
\[
 \frac{\operatorname{Nm}_a(\alpha)}{N^{d_a}}\mathcal D_P(F)
             =(n')^{-d_a}\mathcal D_P(h).
\]
The unit factor has valuation zero.
The assumed chart connection bound gives the conclusion.
For \(P=0\) the value is zero and its valuation is \(+\infty\).
This is the complete calculation in
\texttt{HarmonicValueTranslatedProductValuationBound.term\_bound}.
\end{proof}

\subsection{The actual formal kernel is finite and split}
\label{subsec:ftb-actual-split-kernel}

In an actual good chart, the formal group used in
\S\ref{subsec:ftb-formal-kernel-trace} is the completion at the identity
of the same smooth semi-abelian integral model.
Its scalar kernels on the special fibre are finite: the torus kernel
is a product of \(\mu_N\)'s, and the proper abelian quotient has finite
multiplication kernel.
The exact torus--abelian sequence therefore gives a finite scalar
kernel.  In the formal coordinates its local zero set is supported
at the origin.  Noetherianity implies that the coordinate ideal has a
power contained in the scalar-kernel equations.
Thus this is precisely the finite-height condition used in
Lemma~\ref{lem:ftb-formal-substitution}, including the primes of bad
reduction.

For completeness, the split-kernel assertion in that lemma is
obtained from actual rational torsion, as follows.
Suppose
\[
 \mu'\mu=kM1,\qquad (k,\ell)=1,
\]
at residue characteristic \(\ell\), and every \(M\)-torsion point of
the generic abelian variety is rational over the chart's fraction
field.
Multiplication by \(k\) is an automorphism of the formal group:
its tangent matrix is the unit matrix \(k1\), and the formal inverse
is constructed recursively.
Hence a formal \(\mu\)-kernel point is killed by \(M\).

The formal kernel algebra is finite over \(V\).
A generic algebraic point of it can be taken over a finite extension
DVR, where its coordinates are integral and small.
The formal coordinate point gives compatible points of the same
smooth group scheme modulo all powers of the maximal ideal.
They algebraize over the complete DVR in an affine neighborhood of
the identity: the compatible coordinate ring maps have their limit
in the complete ring and satisfy its finitely many defining equations.
The resulting generic group point is \(M\)-torsion, so by the stated
rational-torsion condition it descends to the original fraction field.
Its chart coordinates descend with it.
Those coordinates are integral over the original \(V\), which is
integrally closed, so they lie in \(V\).
They are still small, since their reductions over the extension are
zero.  This proves the exact \texttt{SplitKer} condition.
It also proves that the formal part of a scalar kernel consists of
its \(\ell\)-primary points; prime-to-\(\ell\) scalar multiplication
is invertible formally.
The source carries out the coordinate/jet and descent argument in
\texttt{SplitKernelFormalGroupOfRationalTorsion}, and derives finite
height in \texttt{FiniteHeightFormalGroupUniformisedModel}.

The rational-torsion condition is reached in the denominator argument
by passing to a finite extension DVR in the actual short-value
reading.  Normalized valuations \(v_{V'}/(te)\) restrict exactly to
\(v_V/e\); no additional denominator constant is introduced by this
extension.

\paragraph{Source ledger and retained boundary.}
The reached-section proof uses the actual N3/PEL expansion
correspondence, integral cubical identifications, primitive theta
generator and period-pairing valuation.
The coefficient calculations above are
\texttt{GeneratorSectionExpansionNuNonneg},
\texttt{NormalisedPullbackSectionNuNonneg},
\texttt{SymmetricTorsionTranslateSectionNuNonneg}, and
\texttt{ReachedSectionRegularOnIntegralModel}.
The exact translation shift and periodic argument are
\texttt{src/ReciprocityTheory/ReductiveGroups/ArithmeticGeometry/DiscreteValuationRings/TWMain.lean}
and \texttt{PeriodicValuation.lean}.
The cubical ratio calculation is
\texttt{src/Development/AbelianVarietyTheory/Proofs/RingsOfIntegers/ValuationThetaRatioEqBilinear.lean}.
The mixed-translate and harmonic statements are in the same
\texttt{RingsOfIntegers/} directory.
None of these deductions by itself constructs the initial integral
uniformization or proves the chart harmonic connection bound
\(\rho\); those are separate producer obligations.
